
\subsection{诺维科夫环与万有诺维科夫环上的弗洛尔上同调}
\label{subsec:5-1-3}

第 2 章在对拉格朗日子流形施加若干限制的条件下，
讨论了以 $\Lambda(L^{(0)},L^{(1)};\ell_0)$
为系数的弗洛尔上同调\liftIndex{I181}的构造。
另一方面，第 3 章对一般的相对 spin 拉格朗日子流形对
构造了滤过 $A_\infty$ 双模，
其系数环是 $\Lambda_{0,\mathrm{nov}}$
或其局部化 $\Lambda_{\mathrm{nov}}$。
若所有障碍均消失，
或者对 $i=0,1$ 两者而言
Maurer--Cartan 模空间 $\mathcal M(L^{(i)})$
均非空，则可形变弗洛尔上边界映射，
从而得到该对子流形的弗洛尔上同调。
本小节说明第 2 章的早期构造
与第 3 章通过 $A_\infty$ 双模所作构造之间的关系。

设 $(L^{(0)},L^{(1)})$ 是横截相交的
拉格朗日子流形对。
首先假定它是相对 spin 对。
此时两个 $L^{(i)}$ 都可定向，
因此交点的马斯洛夫--莫尔斯指标的奇偶性有良定义：
$\mu_{\ell_0}([\ell_p,w])$ 的奇偶性
不依赖于填充道路 $w$ 的选取。
定义 3.7.20 中的模
$C(L^{(1)},L^{(0)};\Lambda_{0,\mathrm{nov}})$，
以及相应的
$C(L^{(1)},L^{(0)};\Lambda_{\mathrm{nov}})$，
分别是在 $\Lambda_{0,\mathrm{nov}}$
与 $\Lambda_{\mathrm{nov}}$ 上
由 $\langle p\rangle$ 自由生成的模。
这里，当 $\mu_{\ell_0}([\ell_p,w])$ 为偶数时，置
\[
\langle p\rangle
=T^{-\mathcal A_{\ell_0}([\ell_p,w])}
 e^{-\mu_{\ell_0}([\ell_p,w])/2}[\ell_p,w];
\]
当 $\mu_{\ell_0}([\ell_p,w])$ 为奇数时，置
\[
\langle p\rangle
=T^{-\mathcal A_{\ell_0}([\ell_p,w])}
 e^{-(\mu_{\ell_0}([\ell_p,w])-1)/2}[\ell_p,w].
\]
于是 $\langle p\rangle$ 的作用量为零，
其次数随该交点的马斯洛夫--莫尔斯指标的
奇偶性分别为 $0$ 或 $1$。
若记
\[
\overline V
=\bigoplus_{k\in\mathbb Z}
 \bigoplus_{p\in L^{(0)}\cap L^{(1)}}
 \mathbb Q e^k\langle p\rangle,
\]
则 $C(L^{(1)},L^{(0)};\Lambda_{0,\mathrm{nov}})$
由 $\overline V$ 作为
$\Lambda_{0,\mathrm{nov}}^{(0)}$-模生成
（参见第 6.3 节）；
$\overline V$ 张成能量为零的子空间。
换言之，这里把 $\langle p\rangle$
视为能量为零的生成元。
由注 3.7.46 可得
\[
C(L^{(1)},L^{(0)};\Lambda_{0,\mathrm{nov}})
=\widehat{\bigoplus}_{[\ell_0]\in
 \pi_0(\Omega(L^{(0)},L^{(1)}))}
C(L^{(1)},L^{(0)};
  \Lambda_{0,\mathrm{nov}};\ell_0),
\]
其中带帽直和表示直和的完备化。

回顾定义 2.4.2 中的弗洛尔上链模。
依定义，有嵌入
\begin{equation}
CF^*(L^{(1)},L^{(0)};\ell_0)
\longrightarrow
C^*(L^{(1)},L^{(0)};\Lambda_{\mathrm{nov}};\ell_0),
\tag{5.1.15.1}\label{eq:5-1-15-1}
\end{equation}
其具体映射是
\begin{equation}
[\ell_p,w]\longmapsto
e^{(\mu([\ell_p,w])-\mu(\langle p\rangle))/2}
T^{\mathcal A_{\ell_0}([\ell_p,w])}
\langle p\rangle.
\tag{5.1.15.2}\label{eq:5-1-15-2}
\end{equation}
该映射与定义 5.1.4 的嵌入
$I_{L^{(0)},L^{(1)};\ell_0}:
\Lambda(L^{(0)},L^{(1)};\ell_0)
\to\Lambda_{\mathrm{nov}}$ 相容。

如果诺维科夫环的底环是域，
则其零次部分
$\Lambda^{(0)}(L^{(0)},L^{(1)};\ell_0)$
也是域，即任意非零元均可逆（参见 [HoSa95]）。
万有诺维科夫环的零次部分
$\Lambda_{\mathrm{nov}}^{(0)}$ 也是域。
由定义 5.1.4 的
$I_{L^{(0)},L^{(1)};\ell_0}$ 可知，
$\Lambda^{(0)}(L^{(0)},L^{(1)};\ell_0)$
嵌入 $\Lambda_{\mathrm{nov}}^{(0)}$，
后者是前者的域扩张。
特别地，$\Lambda_{\mathrm{nov}}^{(0)}$
作为 $\Lambda^{(0)}(L^{(0)},L^{(1)};\ell_0)$-模
是平坦的，上同调的计算因而与系数扩张交换。
另须注意，$\Lambda(L^{(0)},L^{(1)};\ell_0)$
是环
$\bigoplus_k\Lambda^{(0)}(L^{(0)},L^{(1)};\ell_0)e^k$
相对于能量滤过的完备化；
此处 $e$ 不改变能量，只把次数平移 $2$。

现在考虑 $L^{(i)}$ 均为
极小马斯洛夫数至少为 $3$ 的
单调拉格朗日子流形的情形，
如定理 2.4.42 所述。
若假定它们可定向且相对 spin，取
$R=\mathbb Z$；在一般情形下
（不要求它们可定向）取 $R=\mathbb Z_2$。
考虑第 2.3.3、2.4.1 小节定义的模
$CF^*(L^{(1)},L^{(0)};\ell_0)$
及上边界算子
\[
\delta_0:CF^*(L^{(1)},L^{(0)};\ell_0)
\longrightarrow
CF^{*+1}(L^{(1)},L^{(0)};\ell_0).
\]
由定理 2.4.42，在此情形下
$\delta_0\circ\delta_0=0$。
因此弗洛尔上同调
\[
HF(L^{(1)},L^{(0)};\ell_0)
=\operatorname{Ker}\delta_0/\operatorname{Im}\delta_0
\]
有定义，且是
$\Lambda(L^{(0)},L^{(1)};\ell_0)$-模。

若 $L^{(i)}$ 不可定向，
交点的奇偶性未必有良定义：
$\mu_{\ell_0}([\ell_p,w])$
的奇偶性可能依赖于连接 $\ell_0$ 与 $\ell_p$
的填充道路 $w$。
这时稍作修改，
把 $\langle p\rangle$ 定义为
能量和次数均为 $0$ 的形式生成元，
并允许 $e$ 的半整数次幂来扩张万有诺维科夫环。
由此得到嵌入
\begin{equation}
CF^*(L^{(1)},L^{(0)};\ell_0)
\longrightarrow
C^*(L^{(1)},L^{(0)};
\Lambda_{\mathrm{nov}}[e^{1/2}];\ell_0),
\tag{5.1.16.1}\label{eq:5-1-16-1}
\end{equation}
其映射为
\begin{equation}
[\ell_p,w]\longmapsto
e^{\mu([\ell_p,w])/2}
T^{\mathcal A_{\ell_0}([\ell_p,w])}
\langle p\rangle.
\tag{5.1.16.2}\label{eq:5-1-16-2}
\end{equation}
注意，在单调情形下，
第 3 章的构造可在
$R=\mathbb Z$ 或 $\mathbb Z_2$ 上进行，
因此弗洛尔上同调也可在这些底环上定义。

\begin{proposition}{5.1.17}\label{prop:5-1-17}
\textup{(1)} 若 $L^{(i)}$ 的极小马斯洛夫数
均至少为 $3$，则存在定界余链
$b^{(i)}\in C^1(L^{(i)};\Lambda_{0,\mathrm{nov}})$，
它们形如
\[
b^{(i)}=\sum_k T^{\lambda_k}e^{n_k}b_k^{(i)},
\qquad \deg b_k^{(i)}<0.
\]
此外，映射~\eqref{eq:5-1-16-2}
诱导同构\translatornote{原书实体页 283
在以下两条跨连通分支的直和式中
把张量积底环写作
$\Lambda(L^{(0)},L^{(1)})$，
省略了分支参数 $\ell_0$；
按定义 2.2.6 及本页证明末的上链映射，
这里逐分支恢复 $;\ell_0$。
见疑误 A-0265。}：
\[
\begin{aligned}
&\bigoplus_{[\ell_0]\in
 \pi_0(\Omega(L^{(0)},L^{(1)}))}
\Bigl(
HF(L^{(1)},L^{(0)};\ell_0)
\otimes_{\Lambda(L^{(0)},L^{(1)};\ell_0)}
\Lambda_{\mathrm{nov}}^{\mathbb Z_2}[e^{1/2}]
\Bigr)\\
&\hspace{2em}\cong
HF((L^{(1)},b^{(1)}),(L^{(0)},b^{(0)});
\Lambda_{\mathrm{nov}}^{\mathbb Z_2})[e^{1/2}].
\end{aligned}
\]
若 $(L^{(1)},L^{(0)})$ 还构成相对 spin 对，
则另有同构
\[
\begin{aligned}
&\bigoplus_{[\ell_0]\in
 \pi_0(\Omega(L^{(0)},L^{(1)}))}
\Bigl(
HF(L^{(1)},L^{(0)};\ell_0)
\otimes_{\Lambda(L^{(0)},L^{(1)};\ell_0)}
\Lambda_{\mathrm{nov}}^{\mathbb Z}
\Bigr)\\
&\hspace{2em}\cong
HF((L^{(1)},b^{(1)}),(L^{(0)},b^{(0)});
\Lambda_{\mathrm{nov}}^{\mathbb Z}).
\end{aligned}
\]

\textup{(2)} 若对某个哈密顿微分同胚
$\psi$ 有 $L^{(1)}=\psi(L^{(0)})$，
且 $L^{(0)}$ 的极小马斯洛夫数至少为 $2$，
则同样的结论成立。
\end{proposition}

\begin{proof}
极小马斯洛夫数至少为 $3$ 意味着，
只要 $\mathcal M_1(L^{(i)};\beta)\ne\varnothing$
且 $\beta\ne0$，就有
\[
\dim\mathcal M_1(L^{(i)};\beta)
\geq\dim L^{(i)}+1.
\]
因此
\[
\mathfrak m_0^{L^{(i)}}(1)
=\sum_k T^{\lambda_k}e^{n_k}
 \mathfrak m_{0,k}^{L^{(i)}}(1),
\]
其中\translatornote{原书实体页 283
的本次数式用 $\mathbb Q$ 系数，
但上一页已把单调情形的范围扩至
$R=\mathbb Z$ 或 $\mathbb Z_2$。
译文保留原式，不把此有理系数公式
静默当成整系数证明；见疑误 A-0266。}
\[
\begin{aligned}
\mathfrak m_{0,k}^{L^{(i)}}(1)
&\in C^{\dim L^{(i)}-
\dim\mathcal M_1(L^{(i)};\beta)}
(L^{(i)};\mathbb Q)\\
&\subset\bigoplus_{j<0}C^j(L^{(i)};\mathbb Q).
\end{aligned}
\]
于是可归纳证明障碍消失，
并证明存在
$b^{(i)}\in C^1(L^{(i)};\Lambda_{0,\mathrm{nov}})$，
满足
\[
b^{(i)}=\sum_k T^{\lambda_k}e^{n_k}b_k^{(i)},
\qquad\deg b_k^{(i)}<0.
\]
利用 $b_k^{(i)}$ 的上同调次数为负这一事实
（即它的同调维数大于 $\dim L^{(i)}$），
由维数计数可证
\[
\delta_{b^{(1)},b^{(0)}}=\delta_{0,0};
\]
这里 $\delta_{b^{(1)},b^{(0)}}$
见式 (3.7.16)。
依定义，式~\eqref{eq:5-1-16-1}--\eqref{eq:5-1-16-2}
给出上链映射，并诱导同构
\[
\begin{aligned}
&CF^*(L^{(1)},L^{(0)};\ell_0)
\otimes_{\Lambda(L^{(0)},L^{(1)};\ell_0)}
\Lambda_{\mathrm{nov}}[e^{1/2}]\\
&\hspace{3em}\longrightarrow
C^*(L^{(1)},L^{(0)};
\Lambda_{\mathrm{nov}}[e^{1/2}];\ell_0).
\end{aligned}
\]
前述平坦性遂完成命题的证明。
第 \textup{(2)} 项的证明类似。
\end{proof}

\begin{remark}{5.1.18}\label{rem:5-1-18}
若 $L^{(0)}$ 与 $L^{(1)}$ 的马斯洛夫类均为零，
便可按 [Sei00] 的意义为它们赋予分次，
从而使二者的每个交点都有典范的分次。
因此，只要障碍消失，
便可定义以 $\Lambda_{\mathrm{nov}}^{(0)}$
为系数的弗洛尔分次上链复形
$C^*(L^{(1)},L^{(0)};
\Lambda_{\mathrm{nov}}^{(0)})$。
类似地，若 $L^{(0)}$ 与 $L^{(1)}$
的极小马斯洛夫数的最大公因数记为 $2N$，
则可在万有诺维科夫环
$\Lambda_{\mathrm{nov}}$ 的子环
$\Lambda_{\mathrm{nov},N}$ 上定义弗洛尔上链复形，
其中
\[
\begin{aligned}
\Lambda_{\mathrm{nov},N}
=\biggl\{\sum_i a_i e^{\mu_i}T^{\lambda_i}
\in\Lambda_{\mathrm{nov}}\ \bigg|\ 
\mu_i\in N\mathbb Z\text{ 对每个 }i\biggr\}.
\end{aligned}
\]
与上文单调情形相似，
在一般无障碍情形下也可建立同构
\[
\begin{aligned}
&\bigoplus_{[\ell_0]\in
 \pi_0(\Omega(L^{(0)},L^{(1)}))}
\Bigl(
HF(L^{(1)},L^{(0)};\ell_0)
\otimes_{\Lambda(L^{(0)},L^{(1)};\ell_0)}
\Lambda_{\mathrm{nov},N}^{\mathbb Z}
\Bigr)\\
&\hspace{2em}\cong
HF((L^{(1)},0),(L^{(0)},0);
\Lambda_{\mathrm{nov},N}^{\mathbb Z}).
\end{aligned}
\]
第 5.4.6 小节还将进一步讨论此事。
\end{remark}
